\section*{Chapter \ref{ch:g11:os-and-command-line} --- Operating Systems and the Command Line}

\begin{solution}{exo:g11:os-and-command-line:1}
First \code{/home/ada/library/members}: the \code{..} goes up to \code{library}, and the path then goes down
into \code{members}. Then \code{/home/ada/notes}, since two levels up from \code{members} is \code{/home/ada}.
\end{solution}

\begin{solution}{exo:g11:os-and-command-line:2}
\code{/home/ada/library/books/books.csv}: \code{..} leads to \code{/home/ada/library}, and \code{.} stays where it is.
\end{solution}

\begin{solution}{exo:g11:os-and-command-line:3}
\code{rw-} is $4 + 2 = 6$, \code{r--} is $4$, \code{---} is $0$: \code{640}. And $7 = 4 + 2 + 1$, $5 = 4 + 1$,
$4$: \code{rwxr-xr--}.
\end{solution}

\begin{solution}{exo:g11:os-and-command-line:4}
\code{b.txt} and \code{c.txt}, since \code{mv} renamed \code{a.txt}. Then only \code{c.txt}: the second \code{mv}
renames \code{b.txt} to a name that already exists, and the old \code{c.txt}, with the text of \code{a.txt}, is
replaced without a warning.
\end{solution}

\begin{solution}{exo:g11:os-and-command-line:5}
\code{tail -n +2 books.csv | wc -l}, which prints $80$. \code{tail -n +2} prints from the second line on.
\end{solution}

\begin{solution}{exo:g11:os-and-command-line:6}
The \omterm{def:g11:os-and-command-line:shell}{shell} sets up the \omterm{def:g11:os-and-command-line:redirection}{redirection} before \code{cat} starts: \code{> f.txt} empties \code{f.txt} first. \code{cat}
then copies \code{top.txt}, and finds \code{f.txt} empty or, on some systems, refuses with
\code{input file is output file}. Either way, the old lines are lost and \code{f.txt} holds only the top line. Write to another file, then
rename it:

\smallskip\noindent\hspace*{2em}\code{cat top.txt f.txt > new.txt}

\noindent\hspace*{2em}\code{mv new.txt f.txt}
\end{solution}

\begin{solution}{exo:g11:os-and-command-line:7}
\code{ana} is the owner, so the owner's bits \code{rw-} apply: she reads and writes. \code{ben} is in the \omterm{def:g11:os-and-command-line:permission}{group}:
\code{r--}, he reads only. \code{chloe} is among the others: \code{---}, nothing. Readers: \code{ana} and \code{ben};
writer: \code{ana}.
\end{solution}

\begin{solution}{exo:g11:os-and-command-line:8}
\smallskip\noindent\hspace*{2em}\code{cut -d, -f3 books.csv | tail -n +2 | sort | uniq -c | sort -n | tail -3}

\smallskip\noindent It prints \code{6 Elif Ashworth}, \code{6 Quinn Ivymore} and \code{7 Zoe Stoneton}, each count
first. The titles hold no comma, so the third field is always the author.
\end{solution}

\begin{solution}{exo:g11:os-and-command-line:9}
The \omterm{def:g11:os-and-command-line:shell}{shell} splits the line at spaces, so \code{cd} receives two arguments, \code{My} and \code{Notes}, and fails with
\code{too many arguments}. Either quote the name, \code{cd "My Notes"}, or put a backslash before the space, as in
\texttt{cd My\textbackslash{} Notes}; both pass a single argument.
\end{solution}

\begin{solution}{exo:g11:os-and-command-line:10}
The first \code{>>} creates the file, and the second adds to it: \code{g.txt} holds the two lines \code{a} and
\code{b}. With \code{>} the second time, the file is emptied first and holds only \code{b}.
\end{solution}

\begin{solution}{exo:g11:os-and-command-line:11}
\code{chmod g+r,g-w file}. The letters change only the bits named, unlike an octal mode, which sets all nine.
\end{solution}

\begin{solution}{exo:g11:os-and-command-line:12}
Each of the nine bits is $0$ or $1$ independently: $2^9 = 512$ values, the octal numbers \code{000} to \code{777}.
The others may write when the bit of value $2$ in the last digit is $1$; the other eight bits are free, so
$2^8 = 256$ values, half of them.
\end{solution}

\begin{solution}{exo:g11:os-and-command-line:13}
Join the \omterm{def:g11:os-and-command-line:abs-rel}{absolute path} of the \omterm{def:g11:os-and-command-line:abs-rel}{working directory}, a \code{/}, and the \omterm{def:g11:os-and-command-line:abs-rel}{relative path}. Then read the names between the
slashes from left to right, keeping a list of the names kept so far, which starts empty. A \code{.} is skipped. A
\code{..} removes the last name kept, if there is one; at the root there is none, and the root's parent is the root
itself. Any other name is added at the end. The list then never holds \code{.} or \code{..}, and each name kept is a
step down from the one before, so the list, written with slashes, is the \omterm{def:g11:os-and-command-line:abs-rel}{absolute path} of the same file.

Each step reads one name of the joined path, which has finitely many names, so the steps end. (A symbolic link, a
name that points to another place, would make \code{..} go elsewhere; this chapter's file systems have none.)
\end{solution}

\begin{solution}{exo:g11:os-and-command-line:14}
The script is \code{code/grade-11/03-os-and-command-line/exercises/count_loans.sh}:

\smallskip\noindent\hspace*{2em}\texttt{if [ \$\# -ne 1 ]; then}

\noindent\hspace*{4em}\code{echo "usage: bash count_loans.sh MEMBER" >&2}

\noindent\hspace*{4em}\code{exit 1}

\noindent\hspace*{2em}\code{fi}

\noindent\hspace*{2em}\code{cut -d, -f3 loans.csv | tail -n +2 | grep -c -x "$1"}

\smallskip\noindent \texttt{\$\#} is the number of arguments; \code{>&2} sends the message to the standard error, so
that it does not mix with the result; \code{exit 1} stops with a status other than $0$, which means failure.
\code{grep -x} keeps the lines equal to the whole pattern, so that member \code{5} does not count the loans of
\code{55}. For \code{55} the script prints $12$, as the pipeline of the chapter found.
\end{solution}

\begin{solution}{exo:g11:os-and-command-line:15}
\code{sort} compares lines as text, character by character, and \code{9} comes after \code{8}, so \code{9} sorts after
\code{80}. Sorting by numerical value fixes it:

\smallskip\noindent\hspace*{2em}\code{cut -d, -f2 loans.csv | tail -n +2 | sort -n | tail -1}

\smallskip\noindent which prints $80$.
\end{solution}

\begin{solution}{pb:g11:os-and-command-line:1}
\textbf{1.} The owner is \code{ben}: his bits are the owner's. \code{ana} is in the \omterm{def:g11:os-and-command-line:permission}{group} \code{teachers} and is not
the owner: the \omterm{def:g11:os-and-command-line:permission}{group}'s bits. \code{chloe} and \code{dev} are neither: the others' bits.

\textbf{2.} The owner may also be in the \omterm{def:g11:os-and-command-line:permission}{group}. The \omterm{def:g11:os-and-command-line:kernel}{kernel} uses the first class that fits, so an owner gets the
owner's bits even when the \omterm{def:g11:os-and-command-line:permission}{group}'s would give more.

\textbf{3.} She cannot change it: the \omterm{def:g11:os-and-command-line:permission}{group}'s bits \code{r--} apply to her. She can remove it: removing a name changes
the directory, not the file, and her bits on \code{/school}, \code{rwx}, let her.

\textbf{4.} That the \omterm{def:g11:os-and-command-line:permission}{user} has the execute bit of every directory above the file: without it, the file cannot be
reached, whatever its own bits.

\textbf{5.} \code{drwxr-xr-x}, \code{-rw-rw-rw-}, \code{-rw-r--r--}, \code{drwxr-xr-x}, \code{-rw-r--r--},
\code{-rw-r--r--}, in the order of the table.

\textbf{6.} Mode \code{666} gives the others \code{rw-}, and \code{/school} lets everyone enter: any \omterm{def:g11:os-and-command-line:permission}{user} may write
the timetable.

\textbf{7.} \code{/school/marks} has mode \code{755}, so the others may enter it, and \code{g11.csv} has mode
\code{644}, so the others may read it.

\textbf{8.} The student can empty it or write false times in it; everyone would read the false timetable.

\textbf{9.} $r$ lets one list the names inside; $w$ lets one create, remove or rename names inside; $x$ lets one enter
the directory and reach the files inside by their names.

\textbf{10.} With \code{774} the others have \code{r--}: \code{chloe} could list the names, but could not read
\code{g11.csv}, since without $x$ she cannot reach it.

\textbf{11.} With \code{771} the others have \code{--x}: she could not list the names, but she could reach
\code{g11.csv} by its name and, with its mode \code{644}, read it. Hiding a name is not a protection.

\textbf{12.} The students cannot enter \code{/school/marks} any more, so the others' bits of \code{g11.csv} never apply
to anyone: their $2^3 = 8$ values all meet the policy. The owner's bits must be \code{rw-} and the \omterm{def:g11:os-and-command-line:permission}{group}'s
\code{rw-}; the search keeps the one with the fewest bits, \code{660}.

\textbf{13.} Whether a \omterm{def:g11:os-and-command-line:permission}{user} reaches a file depends on the modes of the directories above it, so they must be final
before the file's modes are tried. A directory's path is a beginning of its files' paths, so sorting by length puts it
first.

\textbf{14.} With \code{664}, \code{ana} (owner) has \code{rw-}, \code{ben} (\omterm{def:g11:os-and-command-line:permission}{group}) \code{rw-}, \code{chloe} and
\code{dev} (others) \code{r--}, and all four enter \code{/school}: everyone reads, only teachers write, nobody
executes. Each of its five bits is needed by someone: \code{ana} can get $r$ and $w$ only from the owner's bits,
\code{ben} only from the \omterm{def:g11:os-and-command-line:permission}{group}'s, the students $r$ only from the others'. So every mode that meets the policy has
these five bits, and none may have more: $x$ is forbidden to all, $w$ to the others. \code{664} is the only one.

\textbf{15.} The \omterm{def:g11:os-and-command-line:shell}{commands}, in octal or in letters:
\begin{itemize}
  \item \code{/school/marks}: \code{chmod 770}, or \code{chmod g+w,o-rx};
  \item \code{/school/exam.txt}: \code{chmod 600}, or \code{chmod go-r};
  \item \code{/school/timetable.txt}: \code{chmod 664}, or \code{chmod o-w};
  \item \code{/school/marks/g11.csv}: \code{chmod 660}, or \code{chmod g+w,o-r};
  \item \code{/school} and \code{/school/news.txt} keep \code{755} and \code{644}.
\end{itemize}

\textbf{16.} $755 \to 770$: the \omterm{def:g11:os-and-command-line:permission}{group} gains $w$, the others lose $r$ and $x$, $3$ bits. $644 \to 600$: $2$ bits.
$666 \to 664$: $1$ bit. $644 \to 660$: $2$ bits. In all $3 + 2 + 1 + 2 = 8$, as printed.

\textbf{17.} \code{ana} is not the owner of the exam but is in its \omterm{def:g11:os-and-command-line:permission}{group}, so the \omterm{def:g11:os-and-command-line:permission}{group}'s bits apply to her: they must
be \code{---} for her not to read it.

\textbf{18.} Adding her to the \omterm{def:g11:os-and-command-line:permission}{group} \code{teachers}, which the administrator does once: every right teachers have
comes from the \omterm{def:g11:os-and-command-line:permission}{group}'s bits.

\textbf{19.} A right that nobody needs is never noticed until someone misuses it, and the harm is done before it is
closed. With the fewest rights that work, a mistake or a stolen account can do the least.

\textbf{20.} The policy is met by setting \code{/school/marks} to \code{770}, \code{exam.txt} to \code{600},
\code{timetable.txt} to \code{664} and \code{g11.csv} to \code{660}, keeping \code{755} and \code{644} for
\code{/school} and \code{news.txt}: \textbf{$8$ \omterm{def:g11:os-and-command-line:permission}{permission} bits change}. Every mode is the only one that meets the
policy, except that of \code{g11.csv}, where $8$ do and \code{660} has the fewest bits.
\end{solution}
